%% !TeX program = pdflatex
%% !TeX encoding = UTF-8
%% !TeX spellcheck = pt_BR
%%
%% ficha-exemplo.tex -- PROVA de que \catalogcard poe a ficha catalografica no
%% lugar da moldura "Espaco reservado a ficha catalografica", no PE da folha
%% adicional. A ficha de exemplo e a fichacatalografica.pdf, ao lado deste
%% arquivo (uma pagina A4 qualquer serve de prova).
%%
%% Este exemplo fica em src/ (versionado, para nao envelhecer em silencio), mas
%% NAO vai para as distribuicoes: e prova de funcionamento, nao material do
%% aluno. Para compilar, tenha a classe ao alcance (ex.: copie este .tex e a
%% fichacatalografica.pdf para build/, ou instale a classe) e rode:
%%     pdflatex ficha-exemplo
%% Sem \catalogcard a mesma folha sai com a moldura vazia; com ele, a ficha
%% fornecida ocupa aquele lugar.
\documentclass[dsc]{ufrj}
\usepackage{ufrj-coppe}
\begin{document}
  \title{Prova da ficha catalográfica}
  \foreigntitle{Catalog card placement proof}
  \author{Nome}{Sobrenome}
  \advisor{Nome do Orientador}{Sobrenome}{D.Sc.}{UFRJ}
  \examiner{Nome do Orientador Sobrenome}{D.Sc.}{UFRJ}
  \department{PESC}
  \date{11}{2026}
  \approvaldate{15 de novembro de 2026}
  \keyword{ficha catalográfica}
  \foreignkeyword{catalog card}
  % A PROVA: com \catalogcard, a ficha fornecida substitui a moldura vazia e sai
  % no pe da folha adicional. Comente esta linha e a moldura "Espaço reservado"
  % volta ao mesmo lugar.
  \catalogcard{fichacatalografica.pdf}
  \maketitle
  \frontmatter
  \begin{abstract}
  Documento mínimo para provar a colocação da ficha catalográfica.
  \end{abstract}
  \begin{foreignabstract}
  Minimal document to prove the catalog-card placement.
  \end{foreignabstract}
  \tableofcontents
  \mainmatter
  \chapter{Introdução}
  O corpo existe só para o documento compilar; o que esta prova mostra está na
  folha adicional, antes desta página.
\end{document}
